% Abstract — Manual FATEC 4.1.11 (versão do resumo em língua estrangeira).
% No original, o texto do abstract é apresentado em itálico.
\begin{resumo}[ABSTRACT]
\begin{otherlanguage}{english}
\itshape
The recent Brazilian economic scenario, characterized by macroeconomic
instability, high inflation, and a consequent increase in the Selic interest rate,
poses significant challenges for investors. Financial market volatility,
intensified by internal and external factors, hinders investment decision-making,
especially in assets such as stocks. Given this context and the increasing
digitalization of the financial sector, algorithmic trading presents itself as a promising
alternative. This paper presents the development of a computational tool, using
the MQL5 language on the MetaTrader 5 platform, for analyzing trend and volatility
patterns in the Brazilian stock market, focusing on the PETR4 asset. The tool integrates
statistical technical analyses, combining the use of Moving Averages for
trend identification with the ARIMA model for forecasting price time series and the GARCH
model for modeling the conditional volatility of ARIMA residuals. The central
premise of this work is that this integrated approach allows the generation of trading
signals with greater effectiveness (potential accuracy and risk suitability)
compared to the isolated or incomplete use of these techniques. The study explored
models with updated data, providing analyses that can support investment
strategies adapted to the high-interest-rate scenario, and investigated the
robustness of the combined ARIMA and GARCH models against market fluctuations, contributing to
knowledge in the field and offering a practical alternative for investors.

\vspace{\onelineskip}
\noindent\textbf{Keywords:} Algorithmic Trading. Statistical Models. ARIMA. GARCH.
Stock Market. Volatility. MetaTrader 5.
\end{otherlanguage}
\end{resumo}
